\section*{Chapitre 12 - Opérations sur les fractions}
\addcontentsline{toc}{section}{Chapitre 12 - Opérations sur les fractions}

\subsection*{I. Addition et soustraction de fractions}
\addcontentsline{toc}{subsection}{I. Addition et soustraction de fractions}

\vspace{0.3cm}
\begin{Propriété}
Pour additionner ou soustraire des fractions de même dénominateur (non nul), on ajoute ou soustrait les numérateurs et on conserve le même dénominateur. \\

Quels que soient les nombres $a$, $b$, et $c$ avec $c\neq 0$ : \\
\begin{center}
$ \dfrac{a}{c}+ \dfrac{b}{c}= \dfrac{a+b}{c}$ et $ \dfrac{a}{c} - \dfrac{b}{c} = \dfrac{a-b}{c} $
\end{center}
\end{Propriété}


\begin{Exemples}
\begin{multicols}{2}

$$\begin{aligned}\dfrac{8}{9} - \dfrac{1}{9} &= \ldots \ldots\\
	&=\ldots\ldots
	\end{aligned}$$

\columnbreak

$$\begin{aligned}\dfrac{7}{8} + \dfrac{3}{8}&= \ldots\ldots\\
	&=\ldots\ldots\\
	&=\ldots\ldots
	\end{aligned}$$

\end{multicols}
\end{Exemples}


\begin{Méthode}
Pour additionner ou soustraire des fractions de dénominateurs différents, on les exprime d'abord avec le même dénominateur, puis on applique la propriété précédente.
\end{Méthode}

\begin{Exemples}
{
  \setlength{\columnseprule}{0.4pt} % épaisseur du trait
  \renewcommand{\columnseprulecolor}{\color{Couleur10}}
\begin{multicols}{3}

\begin{align*}
\dfrac{24}{32} + \dfrac{7}{8} &= \ldots\ldots\ldots \\
&=\ldots\ldots \\[0.7em]
&=\ldots\ldots \\[0.7em]
&= \ldots\ldots
\end{align*}


\columnbreak


\begin{align*}
\dfrac{2}{3} - \dfrac{4}{7} &= \ldots\ldots\ldots \\
&= \ldots\ldots \\[0.7em]
&= \ldots\ldots
\end{align*}

\columnbreak


\begin{align*}
4 - \dfrac{7}{9} &= \ldots\ldots \\
&= \ldots\ldots
\end{align*}

\end{multicols}
}
\end{Exemples}


\vspace{0.3cm}
\begin{Remarque}
On dit que l'on \textbf{réduit} deux fractions au même dénominateur lorsqu'on les écrit avec un dénominateur commun.
\end{Remarque}


\newpage
\subsection*{II. Fraction d'un nombre}
\addcontentsline{toc}{subsection}{II. Fraction d'un nombre}

\begin{Propriété}[c]
Pour calculer une fraction d'un nombre, on multiplie cette fraction par ce nombre.
\end{Propriété}

\begin{Méthode}
Calculer $\dfrac{7}{5}$ de 80 revient à calculer $\dfrac{7}{5} \times 80$.\\
Si c'est possible on écrit la fraction sous forme décimale.\\
On a ainsi $\dfrac{7}{5} \times 80 = (7 \div 5) \times 80 = 1,4 \times 80 = 112$.
\end{Méthode}

\begin{Méthode}
On peut aussi interpréter $\dfrac{7}{5}$ de 80 comme 7 fois $ \dfrac{1}{5} $ de 80, c'est-à-dire $7 \times \dfrac{80}{5}$.\\
Dans ce cas, on peut d'abord effectuer la division de 80 par 5, puis multiplier par 7.\\
On a ainsi $\dfrac{7}{5} \times 80 = 7 \times (80 \div 5) = 7 \times 16 = 112$.
\end{Méthode}

\begin{Méthode}
On peut également voir $\dfrac{7}{5} \times 80$ comme $\dfrac{7 \times 80}{5} = \dfrac{560}{5} = 112$.\\
Dans ce cas, on effectue d'abord le produit de 7 par 80, puis on divise par 5.\\
On a ainsi $\dfrac{7}{5} \times 80 = (7 \times 80) \div 5 = 560 \div 5 = 112$.
\end{Méthode}

\begin{Remarque}
Ces trois méthodes donnent le même résultat. Le choix de la méthode ne changera donc pas le résultat sauf si les nombres obtenus dans les calculs intermédiaires sont arrondis. On choisit donc la méthode la plus adaptée à la situation.
\end{Remarque}




